% modeling/sparse-DE.tex
\chapter{Dünn besetztes Alignment-Modell}

\section*{Motivation}

Dieses Kapitel beantwortet, wie kontinuierliche Krümmungstheorie in eine
kompakte Engineering-Darstellung überführt wird, die rekonstruierbar,
interpretierbar und operativ nutzbar bleibt. Es begründet das dünn besetzte
Alignment-Modell, klärt die Verantwortung benachbarter Modellierungsebenen und
bereitet die folgenden Kapitel zu Übergängen und Klassifikation vor.

Die Leitfrage lautet: Welche Information muss persistent sein, und welche
sollte zur Laufzeit rekonstruiert werden?

\begin{figure}[htbp]
\centering
% figures/modeling/sparse_model_structure-DE.tex
\begin{tikzpicture}[
  node distance=0.65cm,
  elem/.style={draw, rounded corners, align=center, minimum width=2.25cm, minimum height=0.9cm},
  arrow/.style={->, thick}
]
\node[elem, fill=blue!8] (f1) {Fest\\\small $G/R$};
\node[elem, fill=orange!12, right=of f1] (t1) {Übergang\\\small $\kappa(s)$};
\node[elem, fill=blue!8, right=of t1] (f2) {Fest\\\small $G/R$};
\node[elem, fill=orange!12, right=of f2] (t2) {Übergang\\\small $\kappa(s)$};
\node[elem, fill=blue!8, right=of t2] (f3) {Fest\\\small $G/R$};
\draw[arrow] (f1) -- (t1); \draw[arrow] (t1) -- (f2);
\draw[arrow] (f2) -- (t2); \draw[arrow] (t2) -- (f3);
\node[draw, rounded corners, fill=green!10, align=center, below=1.2cm of f2, minimum width=6.2cm]
{dünn besetztes Alignment = kompaktes Krümmungsprogramm};
\end{tikzpicture}

\caption{Dünn besetztes Alignment als wechselnde Folge fester und Übergangselemente.}
\label{fig:sparse-model-structure}
\end{figure}

Die erste Abbildung liefert die strukturelle Antwort: Ein dünn besetztes
Alignment ist eine geordnete semantische Folge und kein dichter geometrischer
Schnappschuss.

Ein kontinuierliches Krümmungsgesetz ist mathematisch natürlich, aber für sich
allein keine operative Darstellung für Bearbeitung, Austausch, Validierung
oder Rekonstruktion aus heterogenen Engineering-Daten. Deshalb wird ein dünn
besetztes Modell als strukturierte, stückweise definierte Darstellung eines
Alignments eingeführt.

Das dünn besetzte Modell verneint kontinuierliche Geometrie nicht. Es dient
vielmehr als Zwischenschicht zwischen kontinuierlicher Krümmungstheorie und
praktischen Aufgaben wie Import, Berechnung, Optimierung und
Interoperabilität.

Dichte geometrische Darstellungen wie Punktwolken oder abgetastete Polylinien
bewahren Positionsinformation, zerstören aber weitgehend die strukturelle
Semantik des Alignments selbst.

Die zweite Frage betrifft die Verantwortung: Wie bleiben persistente Daten,
strukturelle Alignment-Semantik und nachgelagerte Engineering Objects getrennt?

\begin{figure}[htbp]
\centering
% figures/modeling/support_alignmentdata_sparse_spotobject-DE.tex
\begin{tikzpicture}[
  node distance=1.25cm,
  box/.style={draw, rounded corners, minimum width=4.2cm, minimum height=0.85cm, align=center},
  arrow/.style={->, thick}
]
  \node[box] (d) {AlignmentData};
  \node[box, below=of d] (s) {SparseAlignment};
  \node[box, below=of s] (o) {SpotObject};
  \draw[arrow] (d) -- (s); \draw[arrow] (s) -- (o);
  \draw[thin] ($(d.center)+(-2.35,0.6)$) rectangle ($(o.center)+(2.35,-0.6)$);
\end{tikzpicture}

\caption{Unterstützende Verantwortungssicht: AlignmentData, SparseAlignment und SpotObject bilden verschiedene Verantwortungsebenen.}
\label{fig:support-alignmentdata-sparse-spotobject}
\end{figure}

Diese Unterscheidung verhindert semantische Drift zwischen Speicherung,
kanonischer Alignment-Semantik und anwendungsspezifischen Laufzeitprodukten.

\section{Grundidee}

\begin{definition}[Dünn besetztes Alignment-Modell]
Ein dünn besetztes Alignment-Modell ist eine strukturierte longitudinale
Darstellung aus synchronisierten Alignment-Bändern mit hinreichender
Rekonstruktionsinformation.
\end{definition}

Das Adjektiv \emph{dünn besetzt} bedeutet hier, dass das Alignment nicht als
dichte Punktwolke, sondern als kompakte Folge semantisch bedeutsamer Elemente
gespeichert wird.

In AIM wird das dünn besetzte Modell nicht als statische 3D-Geometrie, sondern
als persistenter struktureller Kern interpretiert, aus dem Geometrie und
Dynamik abgeleitet werden.

\section{Elementtypen}

\begin{definition}[Alignment-Element]
Ein Alignment-Element ist ein Segment eines Alignments mit vorgegebenem
lokalem Krümmungsverhalten auf einem Parameterintervall.
\end{definition}

\begin{definition}[Festes Element]
Ein festes Element ist ein Alignment-Element mit konstanter Krümmung
\[
\kappaf(s)=\kappaf_0.
\]
Typische Sonderfälle sind Geraden mit \(\kappaf_0=0\) und Kreisbögen mit
\(\kappaf_0\neq 0\).
\end{definition}

\begin{definition}[Übergangselement]
Ein Übergangselement ist ein Alignment-Element mit nicht konstanter Krümmung,
also
\[
\kappaf(s)=f(s)
\]
für ein geeignetes Krümmungsgesetz \(f\).
\end{definition}

In AIM werden Übergangselemente primär durch ihre Krümmungsentwicklung und
nicht durch ihre ausdrückliche geometrische Form interpretiert.

Auf dieser Abstraktionsebene erfolgt noch keine Festlegung auf eine bestimmte
Familie von Übergangsgesetzen. Dies schließt klassische Klothoiden ebenso ein
wie polynomiale, tabellenbasierte oder hybride Konstruktionen.

\TODO{Formale Klassifikation der Übergangsfamilien.}

\section{Minimale Parametrisierung}

\begin{definition}[Bogenlänge eines Elements]
Jedem Element \(e_i\) ist eine Bogenlänge
\[
L_i > 0
\]
zugeordnet.
\end{definition}

\begin{definition}[Lokales Krümmungsgesetz]
Jedem Element \(e_i\) ist ein Krümmungsgesetz
\[
\kappaf_i : [0,L_i] \to \R
\]
zugeordnet.
\end{definition}

Bei festen Elementen ist \(\kappaf_i\) konstant; bei Übergangselementen
verändert es sich entlang des lokalen Parameters.

\begin{definition}[Dünn besetztes Alignment]
Ein dünn besetztes Alignment ist ein geordnetes Tupel
\[
A=(e_1,e_2,\dots,e_n)
\]
mit hinreichenden Anfangsdaten für die geometrische Rekonstruktion.
\end{definition}

\section{Anfangsdaten und Pose}

\begin{definition}[Ebene Anfangspose]
Eine ebene Anfangspose ist ein Tupel
\[
P_0=(x_0,y_0,\theta_0)
\]
mit der Position \((x_0,y_0)\in\R^2\) und der Anfangsrichtung
\(\theta_0\in\R\).
\end{definition}

\begin{corollary}[Rekonstruktionskette]
\label{cor:sparse-reconstruction-chain}
Nach \cref{prop:curvature-driven-reconstruction} wird das ebene Alignment
aus der Anfangspose durch sukzessive Integration des Krümmungsgesetzes
rekonstruiert:
\[
\kappa(s)\rightarrow\theta(s)\rightarrow\gamma(s).
\]
\end{corollary}

Gemeinsam mit der geordneten Elementfolge und den lokalen Krümmungsgesetzen
bestimmt die Anfangspose die Alignment-Geometrie durch sukzessive Integration.

\begin{definition}[Rekonstruierbares dünn besetztes Alignment]
Ein dünn besetztes Alignment heißt rekonstruierbar, wenn seine Elementfolge
und Anfangsdaten genügen, um im Rahmen der vorgesehenen
Modellierungskonventionen ein eindeutiges ebenes Alignment zu bestimmen.
\end{definition}

\section{Dünn besetzte longitudinale Bandstruktur}

\begin{principle}[Prinzip dünn besetzter longitudinaler Bänder]
Eisenbahn-Alignments werden als synchronisierte longitudinale Bänder und nicht
als monolithische geometrische Objekte dargestellt.
\end{principle}

Im AIM-Rahmen wird das persistente Alignment-Modell bewusst dünn besetzt und
in synchronisierte longitudinale Bänder zerlegt.

\subsection{cGeom}

\begin{definition}[Geometrisches Kernband]
Das zentrale geometrische Band wird bezeichnet durch
\[
\mathrm{cGeom}.
\]
Es stellt den krümmungsgetriebenen ebenen Alignment-Kern dar und besteht aus:
\begin{itemize}
\item festen Elementen,
\item Übergangselementen und
\item pose2-basierten Rekonstruktionsankern.
\end{itemize}
\end{definition}

Das cGeom-Band bildet den primären Integrationskern des Alignments.

\subsection{Überhöhungsband}

\begin{definition}[Überhöhungsband]
Überhöhung wird als unabhängiges longitudinales Band dargestellt, das mit
festen Alignment-Elementen korreliert ist.
\end{definition}

Das Überhöhungsband
\begin{itemize}
\item definiert kein unabhängiges Stationierungssystem,
\item speichert Schienenhöhenunterschiede statt Rollwinkel,
\item folgt derselben normierten Krümmungsfamilie wie das zugehörige
cGeom-Element und
\item kann lokale Stationsversätze an Elementgrenzen enthalten.
\end{itemize}

Da die Überhöhungsentwicklung derselben normierten Krümmungsfamilie wie das
zugehörige cGeom-Element folgt, benötigt das Überhöhungsband kein eigenes
geometrisches Entwicklungsmodell.

Der zugehörige Rollwinkel wird abgeleitet durch
\[
\phi(s)=\arctan\!\left(\frac{u(s)}{g}\right),
\]
wobei \(u(s)\) den Schienenhöhenunterschied und \(g\) die Spurweite bezeichnet.

\subsection{Längsprofilband}

\begin{definition}[Längsprofilband]
Die vertikale Geometrie wird unabhängig durch ein Längsprofilband dargestellt.
\end{definition}

cGeom, Überhöhung und Längsprofil werden somit als synchronisierte, aber
verschiedene Bänder mit gemeinsamem Bogenlängengebiet behandelt.

\subsection{Laufzeitzustand}

\begin{proposition}
pose3 wird im dünn besetzten Modell nicht als persistente Wahrheit behandelt.
\end{proposition}

Die direkte Persistierung von pose3 würde Redundanz und
Synchronisationsinstabilitäten zwischen Geometrie-, Überhöhungs- und
Längsprofildarstellung einführen.

Stattdessen wird pose3 dynamisch abgeleitet aus:
\begin{itemize}
\item cGeom,
\item Überhöhung,
\item Längsprofil und
\item Laufzeitbewertungsdiensten.
\end{itemize}

Die primäre Intelligenz des AIM-Rahmens liegt folglich nicht im Speicherformat
selbst, sondern in den mathematischen Operatoren und Bewertungsdiensten, die
auf dem dünn besetzten Modell wirken.

\section{Verkettung}

\begin{definition}[Elementverkettung]
Zwei aufeinanderfolgende Elemente sind verkettet, wenn der Endzustand des
ersten den Anfangszustand des zweiten Elements bereitstellt.
\end{definition}

\begin{definition}[Positionsstetigkeit]
Ein dünn besetztes Alignment ist positionsstetig, wenn aufeinanderfolgende
Elemente in einem gemeinsamen Endpunkt zusammentreffen.
\end{definition}

\begin{definition}[Richtungsstetigkeit]
Ein dünn besetztes Alignment ist richtungsstetig, wenn aufeinanderfolgende
Elemente an ihrer gemeinsamen Grenze dieselbe Tangentenrichtung besitzen.
\end{definition}

In den meisten Engineering-Kontexten sind Positions- und Richtungsstetigkeit
Mindestanforderungen. Höhere Stetigkeitsanforderungen können von den gewählten
Übergangsfamilien und der vorgesehenen Anwendung abhängen.

\TODO{Stetigkeitsklassen wie $C^0$, $C^1$ und krümmungsbezogene Stetigkeitsbedingungen formalisieren.}

\section{Normalform und Wechsel}

In praktischen dünn besetzten Darstellungen ist der Wechsel zwischen festen
und Übergangselementen ein häufiges Muster. Er entspricht weitgehend der
Engineering-Intuition und vielen traditionellen Alignment-Beschreibungen.

\OPEN{Soll der strikte Wechsel zwischen festen und Übergangselementen als mathematische Eigenschaft, bevorzugte Normalform oder lediglich als Implementierungskonvention formuliert werden?}

Es ist plausibel, den Wechsel als Normalform statt als fundamentales Axiom zu
betrachten. Mehrere benachbarte feste Elemente können zusammengeführt werden,
und bestimmte degenerierte Übergangsfälle können in feste Elemente übergehen.

\section{Normierte Darstellung von Elementen}

\begin{definition}[Normierter lokaler Parameter]
Ein normierter lokaler Parameter ist eine Variable
\[
u\in[0,1],
\]
mit der ein Element unabhängig von seiner physischen Länge dargestellt wird.
\end{definition}

\begin{definition}[Normiertes Krümmungsgesetz]
Ein normiertes Krümmungsgesetz ist ein auf dem Einheitsintervall definiertes
Gesetz, typischerweise der Form
\[
\widehat{\kappaf} : [0,1]\to\R.
\]
\end{definition}

Normierung trennt Form von Maßstab. Dieselbe Übergangsfamilie kann dadurch für
verschiedene Längen und Krümmungsbeträge wiederverwendet werden.

Normierte Darstellungen bilden die Grundlage für wiederverwendbare
Übergangsregister, Familienklassifikation und implementierungsunabhängigen
Vergleich von Alignment-Elementen.

\section{Strukturelle Rolle des dünn besetzten Modells}

Das dünn besetzte Modell liegt auf einer begrifflichen Zwischenebene:
\begin{itemize}
\item unterhalb der rein abstrakten krümmungstheoretischen Grundlage,
\item aber oberhalb dichter numerischer Abtastung, importierter Punktlisten
oder formatspezifischer Containerstrukturen.
\end{itemize}

Es ist bewusst weder ein dichtes geometrisches Speicherformat noch eine
direkte Kopie externer Austauschcontainer.

Diese Zwischenstellung macht das dünn besetzte Modell geeignet als
\begin{itemize}
\item interne Modellierungsebene,
\item Rekonstruktionsziel für importierte Daten,
\item Quelle numerischer Bewertung und
\item Verbindung zu Normen und BIM-bezogenen Darstellungen.
\end{itemize}

\section{Bezug zu importierten und standardisierten Daten}

Bestehende Engineering-Daten liegen selten unmittelbar in dünn besetzter Form
vor. Sie werden gewöhnlich durch Containerformate, Vermessungsdaten,
historische Beschreibungen oder abgeleitete Koordinatenfolgen bereitgestellt.

Solche Daten können enthalten:
\begin{itemize}
\item inkonsistente Stationierung,
\item unvollständige Überhöhungsinformation,
\item beschädigte Koordinatenreferenzsysteme,
\item Unstetigkeiten,
\item partielle Geometriefragmente oder
\item nicht synchronisierte Alignment-Komponenten.
\end{itemize}

Das dünn besetzte Modell darf daher nicht mit einem einzelnen Austauschstandard
gleichgesetzt werden. Es ist vielmehr eine kompakte interne Darstellung, die
aus mehreren externen Formaten abgeleitet und in diese zurückabgebildet werden
kann.

\RESEARCH{Dünn besetztes Modell an Alignment-Strukturen von \railML\ anbinden.}
\RESEARCH{Dünn besetztes Modell an Alignment-Darstellungen von \LandXML\ anbinden.}
\OPEN{Eine minimale gemeinsame Abstraktionsebene für Normenabbildungen definieren.}
\OPEN{Den Bezug zwischen topologieorientierten Normen und rein geometrischen Containermodellen klären.}

\section{Bezug zu BIM und BIMalignment}

Ein dünn besetztes Alignment-Modell kann BIM-orientierte Abläufe unterstützen,
sollte aber nicht bloß die Logik von BIM-Containern nachahmen. Seine primäre
Verantwortung ist die mathematisch und operativ kohärente Darstellung der
Alignment-Struktur.

\RESEARCH{Einen BIM-kompatiblen Alignment-Kern formulieren, ohne geometrische Genauigkeit zu verlieren.}
\OPEN{Das dünn besetzte Modell zwischen generischem BIM, Infrastruktur-BIM und ausdrücklichen BIMalignment-Konzepten einordnen.}
\TODO{Ausdrücklich angeben, wo das Manuskript BIMalignment weiterentwickeln und nicht nur BIM-Konformität erfüllen will.}

\section{Zusammenfassung}

\begin{summary}
Das dünn besetzte Alignment-Modell ist eine kompakte, strukturierte Darstellung
eines Alignments durch synchronisierte longitudinale Bänder mit
Rekonstruktionsdaten und Verkettungsregeln.

In AIM wird es nicht als statische 3D-Geometrie, sondern als persistenter
struktureller Kern interpretiert, aus dem Geometrie, Laufzeitzustand und
dynamische Interpretation abgeleitet werden.

Es verbindet kontinuierliche Krümmungstheorie mit praktischen
Engineering-Operationen wie Import, Validierung, Bearbeitung, numerischer
Bewertung, Optimierung und Interoperabilität.
\end{summary}
